\chapter{Conclusion, Future Scope, Limitations}

\section{Conclusion (Current Stage)}
This report has presented the motivation, problem statement, literature
positioning, and complete proposed architecture for the
\ProjectTitle. The system's design -- a state-organized but
mandi-granular forecasting approach, combining a one-time historical
backfill with a continuously updating live pipeline -- was refined
through preliminary analysis of real Agmarknet data rather than fixed
in advance from assumption alone. The project's codebase has been
scaffolded in line with this design and is under active development.

\section{Work Plan and Timeline}
Table~\ref{tab:timeline} presents the project's phase-wise plan and
current status as of \ReviewLabel. Status will be updated at each
subsequent review period.

\begin{longtable}{@{}p{1.2cm}p{7.5cm}p{3cm}@{}}
\toprule
\textbf{Phase} & \textbf{Deliverable} & \textbf{Status} \\
\midrule
\endhead
1 & Supabase connection verified, normalized DB schema (single \texttt{price\_records} fact table, not state-partitioned), Kaggle parquet historical backfill loaded and validated (1,185,657 records, 3 states, 2024-01-01 to 2025-12-29 -- Section~3.2.2) & Done \\
2 & Data-driven top-6-crops-per-state selection (Table~\ref{tab:cropselection}, Chapter~3) & Done \\
3 & Prophet model + naive baseline per (commodity, mandi), MAE comparison & Not Started \\
4 & Sell/Hold decision logic, validated against historical cases & Not Started \\
5 & Retraining pipeline, model registry, state-aggregation logic & Not Started \\
6 & FastAPI backend: per-mandi and state-summary endpoints & Not Started \\
7 & React frontend: Home page & In Progress \\
8 & React frontend: Prediction page (map + drill-down) & Not Started \\
9 & Automated scheduling and deployment & Not Started \\
10 & Polish, documentation, final report & Not Started \\
\bottomrule
\caption{Project phase-wise work plan and status as of \ReviewLabel.}
\label{tab:timeline}
\end{longtable}

\section{Current Limitations}
As the project is at an early implementation stage, the following
limitations apply as of this review and are expected to be resolved in
subsequent phases:
\begin{itemize}
    \item No trained forecasting model exists yet; forecasting accuracy
    cannot be evaluated at this stage (planned for Phase 3).
    \item The live data pipeline and scheduled retraining job are
    currently stub implementations (intentionally not yet scheduled to
    run automatically, to avoid failures on incomplete code).
    \item The frontend currently covers only the Home page structure;
    the interactive India-map prediction page is planned for Phase 8.
    \item Whether multiple reported varieties of the same commodity
    (e.g.\ Onion, Green Chilli) should be collapsed into a single series
    per (commodity, mandi) or modelled separately per variety is not yet
    decided; this affects total model count and will be resolved before
    Phase 3 training begins.
\end{itemize}

\section{Future Scope}
Beyond the current planned scope, the architecture supports several
natural extensions without structural redesign: additional states and
commodities (Section~3.1), explainability layers (e.g.\ SHAP/LIME) for
the forecasting model to support academic evaluation and farmer trust,
SMS/regional-language delivery of Sell/Hold alerts for farmers without
reliable smartphone/internet access, and integration of weather data as
an additional forecasting feature.
